\fichatitulo{22 · RECUPERAÇÃO}{EMNLP · 2020}{Dense Passage Retrieval}
{\small Karpukhin et al. \textit{Dense Passage Retrieval for Open-Domain Question Answering}. EMNLP · 2020. \href{https://aclanthology.org/2020.emnlp-main.550.pdf}{Fonte consultada}.\par}

\campo{Problema e objetivo}
É possível recuperar passagens com representações densas aprendidas?

\campo{Principal contribuição · síntese interpretativa}
Apresenta DPR, com codificadores separados para pergunta e passagem e treinamento supervisionado de suas representações.

\campo{Método / natureza da evidência}
Avaliação em perguntas de domínio aberto, medindo recuperação de passagens e resposta extrativa.

\campo{Resultado ou argumento principal}
O resumo relata ganhos de 9 a 19 pontos percentuais na acurácia de recuperação top-20 em comparação com Lucene-BM25 nos conjuntos avaliados.

\campo{Excertos-chave · seleção literal}
\begin{itemize}
\excerto{1}{a simple dual-encoder framework}{Resumo.}
\excerto{2}{top-20 passage retrieval accuracy}{Resumo.}
\end{itemize}

\campo{Limitações e análise crítica}
Análise da ficha: disponibilidade de pares supervisionados e distribuição dos dados condicionam os ganhos. Encontrar a passagem não assegura fidelidade de uma resposta gerada posteriormente.

\campo{Aproveitamento na dissertação · proposta}
Fundamentação: explicar recuperação densa e separar avaliação do recuperador e do gerador.

\registro{Fonte consultada nesta preparação: Resumo e apresentação da arquitetura. Chave: \texttt{karpukhinEtAl2020dpr}.}
